\documentclass[11pt]{amsart}
\usepackage[T1]{fontenc}
\usepackage[utf8]{inputenc}
\usepackage{amssymb}
\usepackage[margin=1.1in]{geometry}
\usepackage[hidelinks]{hyperref}
\newcommand{\stamptime}{\the\numexpr\time/60\relax:\ifnum\numexpr\time-60*(\time/60)\relax<10 0\fi\the\numexpr\time-60*(\time/60)\relax}
\AddToHook{shipout/foreground}{\put(0.5\paperwidth,-\paperheight+0.55cm){\makebox[0pt]{\tiny\ttfamily unit/doc/\jobname.tex \ --- \ compiled \today, \stamptime}}}
\newtheorem{theorem}{Theorem}[section]
\newtheorem{proposition}[theorem]{Proposition}
\newtheorem{lemma}[theorem]{Lemma}
\newtheorem{corollary}[theorem]{Corollary}
\theoremstyle{definition}
\newtheorem{definition}[theorem]{Definition}
\theoremstyle{remark}
\newtheorem{remark}[theorem]{Remark}
\DeclareMathOperator{\res}{res}
\newcommand{\gcdop}{\operatorname{gcd}}
\newcommand{\cont}{\mathrm c}
\newcommand{\La}{\Lambda}
\newcommand{\N}{\mathbf N}
\newcommand{\Z}{\mathbf Z}
\title{Integral elements from two equations:\\ \v{C}ech cohomology and gcds, constructively}
\date{October 26, 2026}
\begin{document}
\maketitle

\begin{abstract}
Let $k$ be a discrete field, $A=k[a,b]$ and $I=(a,b)$. For two equations
$p=x+af+bg$, $q=y+au+bv$ in $A[x,y]$ we construct, in every $A$-algebra
where $p$ and $q$ vanish, an element $s\equiv1\pmod I$ such that $s$, $sx$
and $sy$ are integral over $A$. The proof is constructive. It computes in an
explicit \v{C}ech model of the projective closure of $p=q=0$, lifts a fixed
cochain by a finite correction whose existence follows from the constructive
noetherianity of $A$, and reads the integral equations on global sections.
Two finiteness statements replace the coherence of higher direct images: one
uses a single element of the elimination ideal, the other uses the gcd of
the homogenised equations. For the latter we prove that $R[X]$ is a gcd
domain when $R$ is one, with Gauss's lemma obtained from the
Dedekind--Mertens lemma, and that a divisor of a homogeneous polynomial is
homogeneous.
\end{abstract}

\section{Statement}

All rings are commutative. A \emph{discrete field} is a ring with $1\ne0$ in
which every element is $0$ or invertible, and this alternative is decidable.

\begin{theorem}\label{thm:main}
Let $k$ be a discrete field, $A=k[a,b]$ and $I=(a,b)$. Let
\[
 P=\mu_P+aF_P+bG_P,\qquad Q=\mu_Q+aF_Q+bG_Q
\]
be homogeneous polynomials in $A[X,Y,Z]$ of degrees $d,e\ge1$, where
$\mu_P=XZ^{d-1}$ and $\mu_Q=YZ^{e-1}$, and let $p=P(x,y,1)$,
$q=Q(x,y,1)$. Let $L$ be an $A$-algebra and $x,y\in L$ with
$p(x,y)=q(x,y)=0$. Then there are $r,t,s\in L$ with
\[
 s=1+ar+bt
\]
such that $s$, $sx$ and $sy$ are roots of monic polynomials in $A[T]$.
\end{theorem}

Taking $L=A[x,y]/(p,q)$ gives the statement for the algebra defined by the
two equations. The existence is not uniform: $s$ is obtained after a
well-founded search, through the constructive noetherianity of $A$.
No regularity of the pair $P,Q$ is assumed; they may have a common factor
(for instance $p=x(1+ax)$, $q=y(1+ax)$, where $Z+aX$ divides $P$ and $Q$).

\subsection*{Plan}
Section~\ref{sec:gcd} proves that $R[X]$ is a gcd domain when $R$ is one,
so that $S=k[a,b][X,Y,Z]$ is a gcd domain.
Section~\ref{sec:laurent} transports the gcd of $P,Q$ to Laurent
polynomials, with homogeneous factors.
Section~\ref{sec:cech} describes the \v{C}ech model.
Section~\ref{sec:lift} lifts the cochain $Z^{-1}$ to a section of
$\mathcal O_W(-1)$, and Section~\ref{sec:integral} produces the integral
equations. Section~\ref{sec:example} follows the construction on an
example.

\section{Gcd domains}\label{sec:gcd}

\begin{definition}
A ring $R$ is a \emph{discrete domain} if one can decide $x=0$ for every
$x$, $1\ne0$, and $x\ne0$, $y\ne0$ imply $xy\ne0$.
For $c,x\in R$, a witness of $c\mid x$ is some $q$ with $x=cq$; divisibility
is never decided. A \emph{gcd} of $x,y$ is an element $d$ with $d\mid x$,
$d\mid y$ and $c\mid d$ for every common divisor $c$ of $x,y$. A \emph{gcd
domain} is a discrete domain together with a function giving a gcd of any
two elements.
\end{definition}

This is the notion of~\cite{MRR}. A gcd is unique up to a unit only; it is data, not a property. Note
$\gcdop(0,0)=0$. A discrete field is a gcd domain: $\gcdop(x,y)=y$ if
$x=0$ and $x$ otherwise.

\begin{lemma}\label{lem:M}
In a gcd domain, $d\mid ux$ and $d\mid uy$ imply $d\mid u\gcdop(x,y)$.
\end{lemma}
\begin{proof}
If $u=0$ there is nothing to prove. Otherwise let $g=\gcdop(ux,uy)$, so
$d\mid g$, and $u\mid g$, say $g=ut$. From $ut\mid ux$ we get $t\mid x$ by
cancelling $u$, and likewise $t\mid y$; hence $t\mid\gcdop(x,y)$ and
$g=ut\mid u\gcdop(x,y)$.
\end{proof}

\begin{corollary}\label{cor:euclid}
If $\gcdop(x,y)\mid1$ and $x\mid yz$, then $x\mid z$.
\end{corollary}
\begin{proof}
$x\mid zx$ and $x\mid zy$, so $x\mid z\gcdop(x,y)\mid z$.
\end{proof}

\subsection*{Content}
For a sequence $e_0,e_1,\dots$ put $\cont_0(e)=0$ and
$\cont_{n+1}(e)=\gcdop(e_n,\cont_n(e))$. By induction, $\cont_n(e)$ divides
$e_i$ for $i<n$, every common divisor of these $e_i$ divides $\cont_n(e)$,
and, by Lemma~\ref{lem:M}, $d\mid ue_i$ for all $i<n$ gives
$d\mid u\cont_n(e)$. For $f\in R[X]$ of degree $<n$ let
$\cont(f)=\cont_n(f_0,f_1,\dots)$. Call $f$ \emph{primitive} if every common
divisor of its coefficients divides $1$.

\begin{lemma}\label{lem:primitive}
(i) If $f$ is primitive and $d\mid uf_i$ for all $i$, then $d\mid u$.
(ii) If $f\ne0$, then $\cont(f)\ne0$ and $f=\cont(f)f^*$ with $f^*$
primitive.
\end{lemma}
\begin{proof}
(i) $d\mid u\cont(f)$ and $\cont(f)\mid1$.
(ii) If $\cont(f)=0$ all coefficients vanish. Divide each coefficient by
$\cont(f)$. If $d$ divides the coefficients of $f^*$, then $d\cont(f)$
divides those of $f$, hence $\cont(f)$, and $d\mid1$ after cancelling
$\cont(f)$.
\end{proof}

\begin{lemma}[Gauss]\label{lem:gauss}
If $E\in R[X]$ is primitive and $d$ divides every coefficient of $Eq$, then
$d$ divides every coefficient of $q$.
\end{lemma}
\begin{proof}
Let $q=\sum_{j\le m}q_jX^j$. The Dedekind--Mertens lemma, in the form of
additive subgroups, which holds over any ring~\cite{mertens}, says that for
every word $w=(i_0,\dots,i_m)$ the product $E_{i_0}\cdots E_{i_m}q_j$ is a
signed sum of products $E_{k_1}\cdots E_{k_m}(Eq)_k$. Since $d$ divides
every $(Eq)_k$, it divides $E_wq_j$ for every word $w$ of length $m+1$.
We now remove the letters of $w$ one at a time. If $d\mid E_wv$ for all
words of length $n+1$, then for each word $w'$ of length $n$ we have
$d\mid E_i(E_{w'}v)$ for all $i$, hence $d\mid E_{w'}v$ by
Lemma~\ref{lem:primitive}(i). For the empty word this gives $d\mid q_j$.
\end{proof}

No prime element and no factorisation are used.

\begin{lemma}\label{lem:divprim}
If $E$ is primitive, $\lambda\ne0$ and $\lambda f=Eq$, then $E\mid f$.
\end{lemma}
\begin{proof}
By Lemma~\ref{lem:gauss}, $\lambda$ divides the coefficients of $q$:
$q=\lambda q'$. Then $\lambda(f-Eq')=0$, so $f=Eq'$.
\end{proof}

\begin{lemma}[elimination]\label{lem:elim}
For $f,g\in R[X]$ one can construct $D\in R[X]$ and $\lambda\ne0$ in $R$
with $D\mid\lambda f$, $D\mid\lambda g$, and such that every common divisor
of $f$ and $g$ divides $D$.
\end{lemma}
\begin{proof}
Induction on $n_f+n_g$, where $f$ (resp.\ $g$) has degree $<n_f$
(resp.\ $<n_g$). If $n_f=0$, then $f=0$: take $D=g$ and $\lambda=1$; the
case $n_g=0$ is symmetric. If $f_{n_f-1}=0$, decrease $n_f$; likewise for
$g$. Otherwise $\alpha=f_{n_f-1}$ and $\beta=g_{n_g-1}$ are non-zero.
Suppose $n_f\ge n_g$ and put $k=n_f-n_g$ and $f'=\beta f-\alpha X^kg$, of
degree $<n_f-1$. Induction gives $D$ and $\lambda'$ for $f',g$. Take
$\lambda=\lambda'\beta\ne0$. Then $\lambda f=\lambda'f'+\alpha X^k\lambda'g$
and $\lambda g=\beta\lambda'g$ are divisible by $D$, and a common divisor
of $f,g$ divides $f'$, hence $D$. If $n_f<n_g$, exchange the roles of $f$
and $g$.
\end{proof}

The field of fractions of $R$ is not used.

\begin{theorem}\label{thm:gcd}
If $R$ is a gcd domain, so is $R[X]$.
\end{theorem}
\begin{proof}
\emph{Discrete domain.} $f=0$ is decidable, since $f$ has finitely many
coefficients. Suppose $f\ne0$, $g\ne0$ and $fg=0$, and let $f_i\ne0$. The
Dedekind--Mertens lemma applied to $fg=0$ gives $f_i^{m+1}g_j=0$ for all
$j$, so $g=0$, a contradiction. (Gauss's lemma alone does not give this.)

\emph{Gcd.} If $f=0$ take $g$, and if $g=0$ take $f$. Otherwise let
$D,\lambda$ be as in Lemma~\ref{lem:elim}. Then $D\ne0$, so
$D=\cont(D)E$ with $E$ primitive, and $E\mid f$, $E\mid g$ by
Lemma~\ref{lem:divprim}. Put $\gamma=\gcdop(\cont(f),\cont(g))$ and
$G=\gamma E$.

$G$ divides $f$: write $f=Eq$. Then $\gamma$ divides the coefficients of
$f$, hence those of $q$ by Lemma~\ref{lem:gauss}. So $q=\gamma q'$ and
$f=Gq'$. The same holds for $g$.

Let $h$ divide $f$ and $g$. Then $h\ne0$ and $h=\cont(h)h^*$. The
coefficients of $f=hr$ are divisible by $\cont(h)$, so
$\cont(h)\mid\cont(f)$; likewise $\cont(h)\mid\cont(g)$, hence
$\cont(h)\mid\gamma$. Moreover $h\mid D=\cont(D)E$, so $h^*\mid\cont(D)E$,
and $h^*\mid E$ by Lemma~\ref{lem:divprim}. Therefore
$h=\cont(h)h^*\mid\gamma E=G$.
\end{proof}

\begin{corollary}\label{cor:S}
For a discrete field $k$, the ring $S=k[a,b,X,Y,Z]$ is a gcd domain.
\end{corollary}

In a gcd domain, let $P\ne0$ and $Q$ be given, $D_0=\gcdop(P,Q)$,
$P=D_0P'$ and $Q=D_0Q'$. Then $D_0\ne0$ and $\gcdop(P',Q')\mid1$: a common
divisor $c$ of $P',Q'$ gives $D_0c\mid D_0$. Hence
\begin{itemize}
\item[(S0)] $D_0f=0$ implies $f=0$;
\item[(S1)] $P'\alpha+Q'\beta=0$ implies $\beta=P'w$ and $\alpha=-Q'w$ for a
  unique $w$.
\end{itemize}
Indeed $P'\mid Q'\beta$ gives $\beta=P'w$ (Corollary~\ref{cor:euclid}); then
$P'(\alpha+Q'w)=0$, and $P'\ne0$.

\section{Homogeneous factors and Laurent polynomials}\label{sec:laurent}

Let $\La=A[X^{\pm1},Y^{\pm1},Z^{\pm1}]$. We represent its elements by their
coefficient functions $\Z^3\to A$ with finite support; the degree of
$m\in\Z^3$ is $m_X+m_Y+m_Z$, and $f$ is homogeneous of degree $k$ if
$f(m)=0$ whenever $\deg m\ne k$. Let $\psi:S\to\La$ send a polynomial to its
coefficient function, extended by $0$ outside $\N^3$. Then $\psi$ is
$A$-linear, injective and multiplicative, and multiplication by $X,Y,Z$
corresponds to the shifts of exponents.

\begin{lemma}[clearing denominators]\label{lem:shift}
For every $f\in\La$ there is $K\ge0$ such that $(XYZ)^Kf=\psi(s)$ for some
$s\in S$: take $K$ at least the sum of the negative parts of the exponents
in a finite list representing $f$.
\end{lemma}

\begin{lemma}\label{lem:reg}
If $s\in S$ is non-zero, then multiplication by $\psi(s)$ is injective on
$\La$.
\end{lemma}
\begin{proof}
If $\psi(s)f=0$, write $(XYZ)^Kf=\psi(s')$. Then $\psi(ss')=0$, so $ss'=0$,
so $s'=0$ since $S$ is a domain, and therefore $f=0$.
\end{proof}

For $x\in S$ let $x_n$ be its homogeneous component of degree $n$. Since $S$
is discrete, $x_n=0$ is decidable, and every $x\ne0$ has a \emph{top degree}
$T$ and a \emph{bottom degree} $B\le T$: $x_T\ne0$, $x_B\ne0$, and
$x_n=0$ for $n>T$ and for $n<B$. These are found by a finite search below a
bound on the degrees of the terms of $x$.

\begin{proposition}\label{prop:homdiv}
Let $x,y\in S$ be non-zero with $xy$ homogeneous of degree $d$. Then $x$
and $y$ are homogeneous.
\end{proposition}
\begin{proof}
Let $T,T'$ be the top degrees of $x,y$, $h=x_T$ and $h'=y_{T'}$. Then
$xy=hh'+h(y-h')+(x-h)y$. The last two terms have no component of degree
$\ge T+T'$, so the component of $xy$ of degree $T+T'$ is $hh'$, which is
non-zero since $S$ is a domain. Hence $T+T'=d$. In the same way the bottom
degrees satisfy $B+B'=d$. Since $B\le T$ and $B'\le T'$, we get $B=T$ and
$B'=T'$, so $x$ and $y$ are homogeneous.
\end{proof}

Apply this to $P=D_0P'$ and $Q=D_0Q'$, with $P,Q$ as in
Theorem~\ref{thm:main} (they are non-zero, since the coefficient of
$\mu_P$ in $P$ is $1$). Then $D_0$, $P'$, $Q'$ are homogeneous, of degrees
$r$, $d'=d-r$ and $e'=e-r$. Through $\psi$ we obtain the following
statements in $\La$:

\begin{proposition}\label{prop:S}
\begin{itemize}
\item[(S0)] If $D_0f=0$ in $\La$, then $f=0$.
\item[(S1)] If $\alpha,\beta\in\La$ are homogeneous of degrees $k+e'$ and
  $k+d'$ with $P'\alpha+Q'\beta=0$, then there is a unique $w\in\La$,
  homogeneous of degree $k$, with $\alpha=-Q'w$ and $\beta=P'w$.
\end{itemize}
\end{proposition}
\begin{proof}
(S0) is Lemma~\ref{lem:reg}. For (S1), multiply $\alpha,\beta$ by the same
$(XYZ)^K$ to land in $\psi(S)$, apply (S1) of Section~\ref{sec:gcd} in $S$,
and divide back by $(XYZ)^K$. Uniqueness follows from Lemma~\ref{lem:reg}
for $P'$. Since $P'$ and $Q'$ are homogeneous, the component of $w$ of
degree $k$ is again a solution, hence equal to $w$.
\end{proof}

\section{The \v{C}ech model}\label{sec:cech}

For a non-empty set $E$ of variables, let $\La_E(k)$ be the elements of
$\La$ of degree $k$ whose exponents of the variables outside $E$ are
$\ge0$. Put
\[
 K^0(k)=\La_X(k)\times\La_Y(k)\times\La_Z(k),\quad
 K^1(k)=\La_{XY}(k)\times\La_{XZ}(k)\times\La_{YZ}(k),\quad
 K^2(k)=\La_{XYZ}(k),
\]
with the usual \v{C}ech differentials $\delta_0,\delta_1$. Two $j$-cochains
are \emph{congruent} if their difference is $Pu+Qv$, componentwise, with
$u,v$ cochains satisfying the same support conditions; we write $N^j(k)$ for
these differences. Then $N^j$ is a submodule stable under $\delta$ and under
multiplication by $X,Y,Z$. A $0$-cochain $c$ is \emph{closed} if
$\delta_0c\in N^1$. The closed cochains of degree $0$ form a ring under the
componentwise product, with unit $(1,1,1)$, and $N^0(0)$ is an ideal of this
ring.

We use three explicit operators, each acting coefficient by coefficient:
the projection $\pi$ of an element of $\La_{XYZ}$ onto its all-negative
monomials, a section $\sigma:K^2\to K^1$ of $\delta_1$, and a contraction
$\tau:K^1\to K^0$. They satisfy, in every degree,
\[
 q=\pi(q)+\delta_1\sigma(q),\qquad
 \delta_0\tau(b)=b\ \text{ if }\delta_1b=0
\]
(the second identity says $H^1(\mathcal O(k))=0$). Moreover
$\pi(\delta_1b)=0$ for every $1$-cochain $b$. Only exponents are compared
in these constructions, never coefficients.

\emph{Restriction to the chart $Z$.} For $f\in\La_Z(k)$ let
$\res(f)=\sum c\,x^{m_X}y^{m_Y}\in L$ over the terms $cX^m$ of $f$. It is
additive and multiplicative, and $\res(P)=p(x,y)=0$, $\res(Q)=0$. Hence the
map $c\mapsto\res(c_Z)$ is a ring map from closed sections of degree $0$ to
$L$ which kills $N^0(0)$. It satisfies
$\res(Xw)=x\res(Zw)$ and $\res(Yw)=y\res(Zw)$ for $w\in K^0(-1)$.

\section{Lifting the cochain $Z^{-1}$}\label{sec:lift}

Let $z_*=(0,0,Z^{-1})\in K^0(-1)$. Then $\delta_0z_*=(0,Z^{-1},Z^{-1})$.

\begin{lemma}[fixed cochain]\label{lem:fixed}
For every $n$ there are $w_n\in I^nK^1(-1)$ and an explicit congruence
$\delta_0z_*\equiv w_n$.
\end{lemma}
\begin{proof}
On the chart $XZ$ put $\theta=\mu_P^{-1}$ and $E=a(-\theta F_P)+b(-\theta
G_P)=1-\theta P$. Then $Z^{-1}=E^nZ^{-1}+Pu_n$ with
$u_n=\theta\sum_{i<n}E^iZ^{-1}$. Similarly, on $YZ$, with $\theta'=\mu_Q^{-1}$
and $E'=1-\theta'Q$. Take $w_n=(0,E^nZ^{-1},E'^nZ^{-1})$, which lies in
$I^nK^1(-1)$ through the expansion of $E^n$ and $E'^n$.
\end{proof}

\begin{lemma}[one element of the elimination ideal]\label{lem:F1}
Let $g=\mu_QP-\mu_PQ$ and $\mathrm{gens}$ the finite set of all-negative
exponents $t$ of degree $-\deg g-1$. Put $m_t=\sigma(gX^t)$. For every $n\ge1$
there are $\rho_n\in A^{\mathrm{gens}}$ with $\rho_{n,t}\in I^{n-1}$, and
$v_n\in I^nK^0(-1)$, such that
\[
 \delta_0z_*\equiv\sum_t\rho_{n,t}m_t+\delta_0v_n .
\]
\end{lemma}
\begin{proof}
$\delta_1w_n=(E'^n-E^n)Z^{-1}$ and $E'-E=\theta P-\theta'Q=g/(\mu_P\mu_Q)$.
Hence $\delta_1w_n=gq_n$ with
$q_n=(\mu_P\mu_Q)^{-1}\sum_{i<n}E'^iE^{n-1-i}Z^{-1}$, of degree $-\deg g-1$,
whose coefficients lie in $I^{n-1}$. Since $\pi(gq_n)=\pi(\delta_1w_n)=0$,
the decomposition $q=\pi(q)+\delta_1\sigma(q)$ applied to $gX^t$ and
$\delta_0\tau=\mathrm{id}$ on cocycles give
\[
 w_n=g\,\sigma(q_n)+\sum_tq_n(t)\,\sigma(gX^t)+\delta_0c_n
\]
with $c_n=\tau(\cdots)\in I^nK^0(-1)$. The first term is congruent to $0$
since $g\in(P,Q)$. Put $\rho_{n,t}=q_n(t)$ and $v_n=c_n$, and use
Lemma~\ref{lem:fixed}.
\end{proof}

Note that $\rho_{n,t}\in I^{n-1}$, so the class of $\delta_0z_*$ in degree
$n$ of the Rees module $\bigoplus_nH^1(I^n\mathcal O_W(-1))$ is generated by
the $m_t$ in Rees degree $1$. No Gröbner basis and no other element of the
elimination ideal is needed.

\begin{lemma}[finite correction]\label{lem:correction}
Let $N\ge1$, $r_i\in I^{N-i}$ for $1\le i<N$, and $v\in I^NK^0(-1)$ with
$\delta_0z_*-\sum_{i<N}r_i\delta_0z_*\equiv\delta_0v$. Then
$h=(1-\sum_ir_i)z_*-v$ is closed and $h-z_*\in IK^0(-1)$.
\end{lemma}
\begin{proof}
$\delta_0h\equiv0$ by linearity, and $r_i\in I$, $v\in IK^0$ since
$N-i\ge1$.
\end{proof}

\begin{proposition}\label{prop:lift}
There is a closed $h\in K^0(-1)$ with $h-z_*\in IK^0(-1)$.
\end{proposition}
\begin{proof}
$A=k[a,b]$ is noetherian in the constructive sense: every sequence of
elements of a finitely generated module, given by codes in a coded
presentation, has a term that is a linear combination of the previous
ones, and this term is reached by a well-founded search. This follows from
the noetherianity of a discrete field, the transfer along $k\to A$, and the
Hilbert basis theorem~\cite{CPersson}. Apply it to the Rees algebra
$\bigoplus_nI^nT^n\subseteq A[T]$, coded inside $A[T]\oplus A[T]^{\mathrm{gens}}$,
and to the sequence $\rho_n$ of Lemma~\ref{lem:F1}. Taking the component of
degree $N$ of the relation found gives $N$ and $r_i\in I^{N-i}$ with
$\rho_N=\sum_{i<N}r_i\rho_i$. Then
\[
 \delta_0z_*-\sum_{i<N}r_i\delta_0z_*\equiv
 \delta_0\Big(v_N-\sum_{i<N}r_iv_i\Big),
\]
with $v=v_N-\sum r_iv_i\in I^NK^0(-1)$. Conclude by
Lemma~\ref{lem:correction}.
\end{proof}

\section{Integral equations}\label{sec:integral}

Let $\delta=d+e-r$, where $r=\deg D_0$, and let $\mathcal G$ be the
all-negative exponents of degree $-\delta$. For $t\in\Z^3$ write $X^t$ for
the corresponding monomial, and put
\[
 c_t=\tau\big(P\,\sigma(-Q'X^t)+Q\,\sigma(P'X^t)\big),\qquad t\in\mathcal G .
\]
For an element $\omega=\sum_t\omega_tX^t$ of degree $-\delta$ with
$\pi(Q'\omega)=\pi(P'\omega)=0$, the cochain $c_\omega=\sum_t\omega_tc_t$ is
closed of degree $0$, with $\delta_0c_\omega=P\sigma(-Q'\omega)+Q\sigma(P'\omega)$
exactly. Indeed $\delta_1$ of the argument of $\tau$ is
$-PQ'\omega+QP'\omega=0$, because $PQ'=QP'=D_0P'Q'$.

\begin{proposition}[(F2)]\label{prop:F2}
For every closed $c\in K^0(0)$ one constructs $\kappa(c)=(a_0,(\omega_t))\in
A^{1+\mathcal G}$ with $c\equiv a_0\cdot1+\sum_t\omega_tc_t$.
\end{proposition}
\begin{proof}
Write $\delta_0c=Pu+Qv$. Then
$0=\delta_1\delta_0c=D_0(P'\delta_1u+Q'\delta_1v)$, so
$P'\delta_1u+Q'\delta_1v=0$ by (S0), and (S1) gives $w$ of degree
$-\delta$ with $\delta_1u=-Q'w$ and $\delta_1v=P'w$. Put
$\omega=\pi(w)$ and $\omega_t=w(t)$. Since $w=\omega+\delta_1\sigma(w)$,
\[
 \delta_1(u+Q'\sigma(w))=-Q'\omega,\qquad \delta_1(v-P'\sigma(w))=P'\omega ,
\]
so $\pi(Q'\omega)=\pi(P'\omega)=0$. Let $u''=u+Q'\sigma(w)-\sigma(-Q'\omega)$
and $v''=v-P'\sigma(w)-\sigma(P'\omega)$. They are cocycles, hence
$u''=\delta_0\tau(u'')$ and $v''=\delta_0\tau(v'')$. Using $QP'=PQ'$ one
finds
\[
 \delta_0c=\delta_0\big(c_\omega+P\tau(u'')+Q\tau(v'')\big),
\]
so $c'=c-c_\omega-P\tau(u'')-Q\tau(v'')$ satisfies $\delta_0c'=0$ exactly.
A $0$-cocycle of degree $0$ satisfying the support conditions is a
constant: its three components agree, and every exponent is $\ge0$ with
degree $0$. Hence $c'=a_0\cdot1$, and $c\equiv a_0+c_\omega$.
\end{proof}

\begin{proof}[Proof of Theorem~\ref{thm:main}]
Let $h$ be given by Proposition~\ref{prop:lift}, and let $c$ be one of the
closed sections $Zh$, $Xh$, $Yh$ of degree $0$. The powers $c^i$ are
closed. Apply the constructive noetherianity of $A$ to the sequence
$\kappa(c^i)$ in $A^{1+\mathcal G}$ (formed with the witnesses of
closedness of the products). It gives $M$ and $\lambda_i\in A$ with
$\kappa(c^M)=\sum_{i<M}\lambda_i\kappa(c^i)$. The map
$(a_0,\omega)\mapsto a_0+c_\omega$ is linear, so
\[
 c^M\equiv\sum_{i<M}\lambda_ic^i\pmod{N^0}.
\]
The restriction is a ring map that kills $N^0$, so $\res(c)$ is a root of
$T^M-\sum_i\lambda_iT^i$. Put $s=\res(Zh)$. Then $\res(Xh)=xs$ and
$\res(Yh)=ys$, so $s$, $sx$, $sy$ are integral over $A$. Finally
$h=(1-\sum r_i)z_*-v$ with $r_i\in I$ and $v\in IK^0(-1)$, and
$\res(Zz_*)=1$, so $s=1-\sum r_i-\res(Zv)\in1+IL$.
\end{proof}

\begin{remark}
The argument uses only the following data about $P$ and $Q$: the shape
$P=\mu_P+aF_P+bG_P$, $Q=\mu_Q+aF_Q+bG_Q$ (for Lemma~\ref{lem:fixed}), and
Proposition~\ref{prop:S} (for Proposition~\ref{prop:F2}). When $D_0$ is a
unit, (S0) is empty and (S1) is the coprimality of $P,Q$; when $D_0$ is not
a unit, as for $p=x(1+ax)$, $q=y(1+ax)$, the quotients $P',Q'$ enter only in
the syzygy, while $P$ and $Q$ still define the congruences and the
sections $c_t$.
\end{remark}

\section{An example}\label{sec:example}

We follow the construction on
\[
 p=x-a+bxy+2bx^2,\qquad q=y-b+ax^2+axy+by^2,
\]
so that $P=XZ-aZ^2+b(XY+2X^2)$ and $Q=YZ+a(X^2+XY)+b(Y^2-Z^2)$, with
$d=e=2$. The numbers below were obtained with an auxiliary computer algebra
script, which checks every identity it uses (in particular
$\delta_1w_n=gq_n$, $\delta_1R=0$ and $\delta_0\tau(R)=R$ for the cochain
$R$ of Lemma~\ref{lem:F1}).

\emph{The gcd.} $\gcdop(P,Q)=1$, so (S0) is empty and (S1) is the
coprimality of $P$ and $Q$. Thus $\delta=4$ and $\mathcal G$ consists of
the three all-negative exponents of degree $-4$; the closed sections are
generated by $1$ and the three $c_t$, in accordance with the Bézout number
$de=4$.

\emph{The generators of (F1).} $g=YZ\,P-XZ\,Q$ has degree $4$, so there are
six generators $m_t$, indexed by the all-negative exponents of degree $-5$,
ordered $(-1,-1,-3)$, $(-1,-2,-2)$, $(-1,-3,-1)$, $(-2,-1,-2)$,
$(-2,-2,-1)$, $(-3,-1,-1)$. The first coordinates are
\begin{align*}
 \rho_1&=(1,0,0,0,0,0),\\
 \rho_2&=(0,b,0,a,0,0),\\
 \rho_3&=(-(a+b)(a+3b),0,b^2,0,ab,a^2),\\
 \rho_4&=(0,-a^3-2a^2b-4ab^2-4b^3,0,-a(a^2+6ab+4b^2),0,0).
\end{align*}
The odd terms live on four generators and the even terms on the other two.

\emph{The relation.} Since $\rho_n$ is homogeneous of degree $n-1$ in
$a,b$, a relation $\rho_N=\sum_{i<N}r_i\rho_i$ with $r_i\in I^{N-i}$ is a
linear system over $\mathbf Q$. There is none for $N\le5$. For $N=6$ one
finds
\[
 \rho_6=r_2\rho_2+r_4\rho_4,\qquad
 r_2=-(a^4+6a^3b+10a^2b^2+32ab^3+9b^4)\in I^4,\quad
 r_4=-2(a+b)(a+3b)\in I^2 ,
\]
and $r_1=r_3=r_5=0$. (The proof does not prescribe which relation is used;
a search through the constructive noetherianity of $A$ may stop at another
index.)

\emph{The element $s$.} The correction is
$h=(1-r_2-r_4)z_*-v$ with $v=v_6-r_2v_2-r_4v_4$, and
$s=\res(Zh)=1-r_2-r_4-\res(Zv)$ is a polynomial in $x,y$ of total degree
$6$, with $52$ terms and coefficients of degree at most $6$ in $a,b$.
Modulo $I^4$,
\[
 s\equiv1+2(a+b)(a+3b)=1-r_4 ,
\]
so $s\equiv1\pmod{I^2}$.

\emph{Integrality.} Over $\mathbf Q(a,b)$ the algebra $A[x,y]/(p,q)$ has
dimension $4$, with basis $1,y,y^2,x$. The characteristic polynomials of
$s$, $sx$ and $sy$ are monic of degree $4$ with coefficients in
$\mathbf Q[a,b]$, as Theorem~\ref{thm:main} predicts.

\emph{Comparison.} An element with the same properties is
$u=(1+ax+by)(1+2bx+by)^2$. The element $s$ produced by the construction is
different: $s-u\notin(p,q)$. Both satisfy equations of degree $4$, but that
of $s$ is much larger; for instance its constant term has degree $24$ in
$a,b$, against $12$ for $u$. The construction is uniform in $p,q$, but it
does not look for a small certificate.

\section{Formalisation}

Everything above is formalised and checked in the proof checker Unit, in
the developments \texttt{sag/gcd} (Sections~\ref{sec:gcd}
and~\ref{sec:laurent}) and \texttt{sag/acl-zmt}
(Sections~\ref{sec:cech}--\ref{sec:integral}). The foundations are those of
the other developments of \texttt{sag}: univalence and replacement are
parameters, function extensionality is derived, and $A[X]$ is a set
quotient of coefficient lists. No other axiom, no choice principle and no
excluded middle is used. The main checked statements are:
\begin{itemize}
\item \texttt{PolyGcd.Over[R,G].domain}: Theorem~\ref{thm:gcd};
  \texttt{Iterate}: Corollary~\ref{cor:S};
\item \texttt{HomDivisor.Divisor}: Proposition~\ref{prop:homdiv};
  \texttt{HomLambda.Data}: Proposition~\ref{prop:S}, with the homogeneous
  factors as lists of terms;
\item \texttt{CechSections.Sec.F2D}: Proposition~\ref{prop:F2};
  \texttt{CechZMTD0}, \texttt{ACLGeneral}: the assembly over a gcd domain
  with a coded noetherian presentation;
\item \texttt{ACLGeneralK.Over[k,\dots].Charts[\dots].certificate}:
  Theorem~\ref{thm:main}, with the charts $\mu_P,\mu_Q$ slightly more
  general than $XZ^{d-1}$, $YZ^{e-1}$;
\item \texttt{ACLFinalGeneral} and \texttt{ACLControlGeneral}: the
  instances $p=x-a+bxy+2bx^2$, $q=y-b+ax^2+axy+by^2$, and
  $p=x(1+ax)$, $q=y(1+ax)$, where all hypotheses are discharged by
  computation.
\end{itemize}
The computations of Section~\ref{sec:example} are not part of the
formalisation; the script that produced them is
\texttt{sag/acl-zmt/acl\_cech\_trace.py}. Divisibility is never decided. The gcd is computed by
Lemma~\ref{lem:elim}, and the degree of a polynomial, being unique, is
extracted from the set quotient. Checking the two developments takes about
seven seconds.

\begin{thebibliography}{9}
\bibitem{CPersson} T.~Coquand and H.~Persson, \emph{Gröbner bases in type
  theory}, in Types for Proofs and Programs, LNCS 1657, Springer, 1999.
\bibitem{mertens} \emph{A direct proof of the Dedekind--Mertens lemma},
  note of April 5, 2006; see also \texttt{unit/doc/dedekind-mertens.tex}.
\bibitem{MRR} R.~Mines, F.~Richman and W.~Ruitenburg, \emph{A Course in
  Constructive Algebra}, Springer, 1988.
\end{thebibliography}

\end{document}
